% Couverture et quatrième de couverture : un aplat rouge, le titre en petit, le
% second mot en très grand, couché le long du bord droit.
\definecolor{toncouv}{HTML}{C33A28}   % cinabre de commutator.io
\definecolor{toncouvencre}{HTML}{5E1A12}

\newcommand{\couverture}{%
  \thispagestyle{empty}
  \begin{tikzpicture}[remember picture, overlay]
    \fill[toncouv] (current page.south west) rectangle (current page.north east);
    % le grand mot, couché, en encre sombre
    \node[rotate=90, anchor=south, inner sep=0pt, text=toncouvencre,
          font=\firathin\fontsize{118}{118}\selectfont]
      at ($(current page.south east)+(-0.9cm,12.4cm)$) {outillées};
    % le titre en petit, et le filet qui le relie au grand mot
    \node[anchor=base west, inner sep=0pt, text=white, font=\sffamily\Large]
      (titre) at ($(current page.north west)+(2.2cm,-6cm)$) {Mathématiques};
    \draw[toncouvencre, line width=1.6pt]
      ($(titre.base east)+(0.35cm,0.12cm)$) -- ($(current page.east|-titre.base)+(-4.6cm,0.12cm)$);
    \node[anchor=north west, inner sep=0pt, text=white!88, font=\sffamily\small,
          text width=8cm, align=left]
      at ($(titre.base west)+(0,-0.6cm)$) {le génie logiciel au service de la recherche};
    % les auteurs en bas à droite
    \node[anchor=south east, inner sep=0pt, text=white, font=\sffamily\normalsize,
          align=left]
      at ($(current page.south east)+(-5.4cm,2.2cm)$)
      {Michel Hua\\[2pt]{\small\color{white!80} avec Claude comme assistant}\\[2pt]
       {\footnotesize\color{white!70} version \versionlivre}};
    % la marque, discrète, en bas à gauche
    \node[anchor=south west, inner sep=0pt, text=white!75, font=\sffamily\scriptsize]
      at ($(current page.south west)+(2.2cm,2.2cm)$) {commutator.io};
  \end{tikzpicture}
  \null\clearpage
}

\newcommand{\quatrieme}{%
  \cleardoublepage
  \thispagestyle{empty}
  \begin{tikzpicture}[remember picture, overlay]
    \fill[toncouv] (current page.south west) rectangle (current page.north east);
    \node[anchor=north west, text width=11.6cm, align=justify, text=white,
          font=\normalsize, inner sep=0pt]
      at ($(current page.north west)+(2.4cm,-3.4cm)$) {%
      {\sffamily\Large Mathématiques outillées}\par\bigskip
      Un assistant de code lit des articles, transcrit des notes manuscrites, écrit du LaTeX
      qui compile et démontre des énoncés en Lean. Il écrit vite, souvent juste, et ses
      erreurs ont la même assurance que le reste.\par\medskip
      Ce livre s'adresse aux mathématiciennes et aux mathématiciens qui veulent s'en servir
      sans perdre la trace de ce qui a été vérifié. Il emprunte au génie logiciel ses outils
      les plus simples : le dépôt git comme cahier de laboratoire, les issues comme carnet de
      recherche, l'intégration continue, Lean et mathlib pour vérifier, une page locale pour
      relire. Il en tire des règles : un en-tête sur chaque fichier, le doute marqué plutôt
      que résolu, aucune nouveauté affirmée sans recherche, rien de publié qu'on n'a pas le
      droit de publier.\par\medskip
      Aucune formation en informatique n'est supposée. Il faut seulement accepter d'ouvrir
      un terminal, et de le faire souvent.};
    \draw[toncouvencre, line width=1.6pt]
      ($(current page.south west)+(2.4cm,3.4cm)$) -- ++(3cm,0);
    \node[anchor=south west, inner sep=0pt, font=\sffamily\footnotesize, text=white!80]
      at ($(current page.south west)+(2.4cm,2.2cm)$) {version \versionlivre\ · 2026};
    \node[anchor=south east, inner sep=0pt, font=\sffamily\scriptsize, text=white!75]
      at ($(current page.south east)+(-2.4cm,2.2cm)$) {commutator.io};
  \end{tikzpicture}
  \null
}
